\section{研究计划与阶段性产出}

\subsection{总体时间安排}


\par 本研究计划于2026年3月4日启动，预计5月11日完成毕业论文初稿，周期约10周。整个周期划分为五个阶段，各阶段任务明确，产出清晰，确保研究进度可控且严格遵循学院关键时间节点。具体安排如下：
\par 第一阶段（3.4-3.23）：开题报告撰写与完善。 这一阶段的核心任务是完成开题报告终稿。我们将根据指导教师反馈完成开题报告的撰写与修改，确保3月23日前提交符合要求的开题报告初稿，并在此基础上进一步完善，形成最终版本。同时，初步梳理数据采集方案，为后续研究奠定基础。产出：开题报告终稿、数据采集初步方案。
\par 第二阶段（3.24-4.6）：数据采集与特征体系构建。 这一阶段的核心任务是将候选地址池扩展至200个以上，并完成特征体系的构建。我们将从Blackbook、Solscan标签、Dune查询和反向追踪四个渠道同步采集地址，通过滑动窗口检测筛选满足连续7天≥1800笔交易的活跃地址。在前期12维特征的基础上，新增优先费用统计、新代币交互比例、程序调用熵等Solana特有指标，最终形成20+维的特征集。产出：addresses.csv（≥200地址）、features\_v2.csv（20+维）、特征分布图与相关性热力图。
\par 第三阶段（4.7-4.20）：监督学习模型构建与调优。 这一阶段的核心任务是完成创新性方法的设计与实现，对应4月20日前的时间节点。我们将构建规模约300-400的标注集，以社区列表中的地址为正样本，从普通用户地址中随机抽取负样本，并进行人工抽样验证。随后训练随机森林、XGBoost、LightGBM三种模型，通过5折交叉验证和网格搜索优化超参数，并与启发式规则基线进行对比。产出：监督学习模型性能对比表、特征重要性排序图、创新性方法说明文档。
\par 第四阶段（4.21-5.4）：无监督聚类与机器人行为分型。 采用高斯混合模型（GMM）和DBSCAN对全部地址进行聚类，探索数据的内在结构。通过分析每个簇的特征均值，尝试解释簇对应的行为类型（狙击型、套利型、做市型等）。若聚类结果揭示出清晰的机器人亚型，将基于聚类结果结合人工验证构建多分类标注集，训练多分类模型。产出：聚类结果与簇分布图、机器人类型划分标准文档、多分类模型初步性能。
\par 第五阶段（5.5-5.11）：毕业论文初稿撰写。 这一阶段的核心任务是完成毕业论文初稿，对应5月11日前的时间节点。集中撰写论文的各个章节：绪论、文献综述、研究方法、实验结果、结论与展望。将前期产出的图表和分析结果整合进论文，形成逻辑自洽的完整论述。产出：毕业论文初稿。

\subsection{预期成果}

\par 本研究预期产出四项成果。

\par 成果一：公开的Solana机器人地址数据集。 该数据集包含200个以上经过活跃度筛选的疑似机器人地址，每个地址附带来源标签、活跃窗口信息以及20余维行为特征。数据集将以CSV格式公开发布，供社区后续研究使用。
\par 成果二：开源的Solana MEV机器人检测pipeline。 整套代码（数据采集、特征计算、模型训练、可视化）将托管于GitHub，提供一键运行脚本和详细的README文档，确保任何研究者都能复现本文结果。
\par 成果三：Solana机器人行为模式与生态影响分析报告。 报告将系统描述Solana机器人的行为特征分布、不同类型机器人的行为分化，以及机器人活动对Meme币生态的影响。